% GPT-6 Astra Pro; effort not exposed; /root/replacement; 22 November 2026.
\section[HD-F1-R20: unrestricted adjacent coupling tangents]{HD-F1-R20: the actual point-action and coupling tangents are disjoint for every adjacent balanced object}
\label{sec:followup-rm-unrestricted-adjacent-tangent}

Retain the genuine RM fourfold and a specified coherent
faithful finite point action. Let an actual gluable object
$E$ with a weakly semiregular lift have precisely the two
nonzero coherent cohomology sheaves $H^0(E)=A$ and
$H^1(E)=T$, and whole character $Ne^{-B}v_q$, $N\ne0$.
The forward invariant cross-Hom group may have arbitrary
dimension. There is no endpoint semistability assumption
and no assertion
that either endpoint has character proportional to that
of $E$. We prove an exact tangent and Euler restriction
at this unrestricted adjacent scope.

Use the actual triangle
\[
       A\xrightarrow{i}E\xrightarrow{p}C=T[-1]
                    \xrightarrow{\delta}A[1]
 \tag{F1R20.1}
\]
and the notation of R19:
\[
 \begin{gathered}
 V=\operatorname{Ext}^2_G(T,A),\quad
 R_A=\operatorname{End}_G(A),\quad R_T=\operatorname{End}_G(T),\quad
 d=\dim R_A+\dim R_T,\\
 d_\delta(u,v)=u[2]\delta-\delta v,\quad
 s=\dim\ker d_\delta,\quad L_\delta=\operatorname{im}d_\delta,\\
 a=\operatorname{ext}^1_G(T,A),\quad
 b=\operatorname{ext}^1_G(A,T),\quad
 h=\operatorname{hom}_G(A,T),\quad c=\chi_G(T,A).
 \end{gathered}
 \tag{F1R20.2}
\]
The exact negative-Ext isomorphism (F1R19.14) shows
that gluability is precisely
$\operatorname{Hom}_G(T,A)=0$. The integer $h$ is
otherwise arbitrary.

\paragraph{1. Both endpoint injections give the stronger disjointness.}
The endomorphism sequence and coupling tangent map
remain the actual maps (F1R18.8)--(F1R18.10), because
their exactness uses only negative Hom groups of sheaves:
\[
 e_0=a+s,\qquad
 J:V/L_\delta\lhook\joinrel\longrightarrow
                       \operatorname{Ext}^1_G(E,E),\qquad
                       J(v)=i[1]vp.
 \tag{F1R20.3}
\]
Suppose $\operatorname{ev}_E^1(x)=J(v)$, where
$x\in HH^1(X)$. Compose first with $i$ and then with
$p[1]$. Since $pi=0$, characteristic naturality gives
\[
           i[1]\operatorname{ev}_A^1(x)=0,\qquad
                    \operatorname{ev}_C^1(x)p=0.
 \tag{F1R20.4}
\]
The two maps detecting these equalities are injective:
\[
 \begin{split}
 \operatorname{Ext}^1_G(A,A)&\longrightarrow
                              \operatorname{Hom}_G(A,E[1]),\\
 \operatorname{Ext}^1_G(C,C)&\longrightarrow
                              \operatorname{Hom}_G(E,C[1]).
 \end{split}
 \tag{F1R20.5}
\]
For the first, the preceding term is
$\operatorname{Hom}_G(A,C)=\operatorname{Ext}^{-1}_G(A,T)=0$.
For the second, the preceding term is
$\operatorname{Hom}_G(A[1],C[1])=\operatorname{Hom}_G(A,C)=0$.
Thus both endpoint characteristic Ext classes vanish.
The shift identification of the second endpoint,
including the graded-centre sign, makes
$\operatorname{ev}_T^1(x)=0$ as well. Taking their
ordinary full traces gives
$c^1_{\operatorname{ch}(A)}(x)=c^1_{\operatorname{ch}(T)}(x)=0$.

Actual $K_0$ additivity along (F1R20.1) gives
$\operatorname{ch}(E)=\operatorname{ch}(A)-\operatorname{ch}(T)$,
so $c^1_{\operatorname{ch}(E)}(x)=0$. The balanced
whole-character map has rank eight on the eight-dimensional
space $HH^1(X)$, hence $x=0$. We have proved the actual
injection
\[
 HH^1(X)\oplus V/L_\delta
       \lhook\joinrel\longrightarrow\operatorname{Ext}^1_G(E,E),
 \qquad(x,[v])\longmapsto\operatorname{ev}_E^1(x)+i[1]vp.
 \tag{F1R20.6}
\]
This proof uses the character of $E$ and both endpoint
maps; no injectivity of an individual endpoint character
map is presumed. Consequently
\[
                            e_1\ge8+\dim V-d+s.
 \tag{F1R20.7}
\]

\paragraph{2. The unrestricted integral Euler inequality.}
Invariant CY4 duality, gluability and full weak
semiregularity give exactly
\[
 \dim V=c+a+b-h,\qquad e_2=18,\qquad
 k:=\chi_G(E,E)=2e_0-2e_1+18\in2\mathbf Z_{>0}.
 \tag{F1R20.8}
\]
Positivity here concerns $k$, whose actual full balanced
character and S-21 Euler map give
$k=4Dq^2N^2/|G|$. We impose no sign on the unrestricted
cross Euler integer $c$. Equations (F1R20.3),
(F1R20.7) and (F1R20.8) give
\[
                            c+b+\frac{k}{2}\le d+1+h.
 \tag{F1R20.9}
\]
There is also a version involving only endpoint self-Euler
numbers. The actual canonical triangle in equivariant
$K_0$ and invariant CY4 duality give
\[
 k=\chi_G(A,A)+\chi_G(T,T)-2c,
 \qquad
 \frac{\chi_G(A,A)+\chi_G(T,T)}2+b\le d+1+h.
 \tag{F1R20.10}
\]
Every term is the Euler pairing of the specified actual
objects and action. No unproved endpoint numerical
character is substituted in (F1R20.9) or (F1R20.10).

\paragraph{3. The minimal three-dimensional endomorphism case.}
For the remaining minimal profile
$(e_0,e_1,e_2)=(3,8,18)$, (F1R20.6) forces
\[
                         V=L_\delta,\qquad a+s=3.
 \tag{F1R20.11}
\]
In particular every actual first-order change of the
coupling class is an endpoint-coordinate change. The
precise action is
\[
 R_A^*\times R_T^*\longrightarrow V,\qquad
           (u,v)\longmapsto u[2]\delta v^{-1},
 \tag{F1R20.12}
\]
whose derivative at $(1,1)$ is the surjective map
$d_\delta$ of (F1R20.2). This is a statement about the
actual coupling and the actual automorphism groups,
not a formal replacement of the triangle.

Let
\[
 S_\delta=\{(u,v)\in R_A\oplus R_T:
                         u[2]\delta=\delta v\}.
 \tag{F1R20.13}
\]
The map from $\operatorname{End}_G(E)$ to its induced
endomorphisms on the two cohomology sheaves is a
surjective unital algebra map to $S_\delta$, with
kernel
\[
 I=\{ifp:f\in\operatorname{Ext}^1_G(T,A)\},\qquad I^2=0.
 \tag{F1R20.14}
\]
Surjectivity and injectivity of the displayed kernel
are the exact sequence (F1R18.8). The square-zero
identity follows directly from
$(ifp)(if'p)=if(pi)f'p=0$.

R07 now makes $\operatorname{End}_G(E)$ one of the two
local three-dimensional algebras there. For negative
$N$, apply that result to the actual shifted lift
$E[1]$, whose self-Ext algebra is identified by the
shift equivalence and whose character has positive rank.
The quotient $S_\delta$ is therefore local. Its identity
has positive dimension, so $0\le a\le2$, and the only
remaining algebra profiles are
\[
 \begin{array}{c|c|c}
 a=\dim I&s=\dim S_\delta&S_\delta\\ \hline
 2&1&\mathbf C\\
 1&2&\mathbf C[\epsilon]/(\epsilon^2)\\
 0&3&\mathbf C[\epsilon]/(\epsilon^3)
       \text{ or }\mathbf C[x,y]/(x,y)^2.
 \end{array}
 \tag{F1R20.15}
\]
In the first line $I$ is the whole two-dimensional
radical, so (F1R20.14) additionally forces
$\operatorname{End}_G(E)=\mathbf C[x,y]/(x,y)^2$.
For the middle line, a two-dimensional unital local
complex algebra has one-dimensional radical, which
is nilpotent and hence square zero, proving the
listed algebra without an extra classification premise.
Finally (F1R20.8) and (F1R20.11) give the exact equality
\[
                              c+b-h=d-3.
 \tag{F1R20.16}
\]

\paragraph{Framework and precise scope.}
The proof uses the same published characteristic/trace
naturality and full HKR action already checked in
R07--R11, and the explicit actual Hom sequences of
R18--R19. Additivity is the class relation in the
actual equivariant $K_0$ of the canonical truncation
triangle. The finite action, its chosen coherence,
and all endpoint and coupling maps are retained.
The original R19 rank-one calculation is not invoked
to prescribe the endpoint characters here.

Equations (F1R20.6), (F1R20.9) and
(F1R20.11)--(F1R20.16) are proved necessary restrictions
for every actual adjacent balanced object in the stated
scope, including larger invariant forward Hom and
nonsemistable endpoints. No geometric realization of
any remaining algebra profile is asserted. These maps
do not prove an unrestricted adjacent-object exclusion
or supply a new weakly semiregular RM seed.
